% !TeX program = xelatex
% ============================================================
%  Causal Inference: What If —— 中文译稿 共享导言区
%  版式特征：外侧宽边距 + 边注（marginalia）、精要框、技术注记框
% ============================================================
\usepackage{amsmath,amssymb,amsthm,mathrsfs,bm}
\usepackage{mathtools}
\usepackage{tikz}
\usetikzlibrary{arrows.meta,positioning,calc,fit,backgrounds,shapes.geometric}
\usepackage{graphicx}
\usepackage{booktabs,array,multirow,longtable}
\usepackage{enumitem}
\usepackage{multicol}
\usepackage{xcolor}
\usepackage[most]{tcolorbox}
\usepackage{caption}
\usepackage{ragged2e}
\usepackage{fancyhdr}
\usepackage{emptypage}
\usepackage{marginnote}
\usepackage{etoolbox}

% ---------- 颜色 ----------
\definecolor{inkblue}{RGB}{30,70,120}
\definecolor{fpbg}{RGB}{235,238,243}   % 精要框底色（原书浅灰蓝）
\definecolor{tpbg}{RGB}{250,250,248}
\definecolor{rule}{RGB}{120,120,120}

% ---------- 页面：外侧留宽边距放边注 ----------
\usepackage[a4paper,inner=2.3cm,outer=5.3cm,top=2.5cm,bottom=2.5cm,
            marginparwidth=3.1cm,marginparsep=0.45cm,
            headsep=0.65cm,footskip=0.95cm]{geometry}
\linespread{1.28}
\let\cleardoublepage\clearpage
\setlength{\parindent}{2em}

% ---------- 标题格式 ----------
\ctexset{
  chapter = {
    format     = \raggedright\heiti\zihao{2},
    name       = {第,章},
    number     = \chinese{chapter},
    aftername  = \hspace{0.55em},
    beforeskip = 0.6cm,
    afterskip  = 0.9cm,
  },
  section = {
    format = \heiti\zihao{3},
    beforeskip = 1.2ex plus .3ex minus .2ex,
    afterskip  = .9ex plus .2ex,
  },
  subsection = {
    format = \heiti\zihao{4},
    beforeskip = 1.0ex plus .3ex minus .2ex,
    afterskip  = .6ex plus .2ex,
  },
  figurename = 图,
  tablename  = 表,
}

\captionsetup{font=small,labelsep=quad}
\captionsetup[figure]{justification=centering}

% ---------- 页眉页脚 ----------
\pagestyle{fancy}
\fancyhf{}
\fancyhead[LE]{\small\kaishu\leftmark}
\fancyhead[RO]{\small\kaishu\rightmark}
\fancyfoot[LE,RO]{\small\thepage}
\renewcommand{\headrulewidth}{0.4pt}
\fancypagestyle{plain}{\fancyhf{}\fancyfoot[LE,RO]{\small\thepage}\renewcommand{\headrulewidth}{0pt}}
\makeatletter
\renewcommand{\chaptermark}[1]{\markboth{\CTEXthechapter\quad #1}{}}
\renewcommand{\sectionmark}[1]{\markright{\thesection\quad #1}}
\makeatother

% ---------- 边注 ----------
\renewcommand*{\marginfont}{\sffamily\footnotesize\raggedright\linespread{1.06}\selectfont}
\newcommand{\mnote}[1]{\marginnote{\textcolor{black!85}{#1}}}
\setlength{\marginparpush}{6pt}

% ---------- 精要框 / 技术注记框 ----------
\newcommand{\FPname}{精要}
\newcommand{\TPname}{技术注记}
\newcounter{fp}[chapter]
\newcounter{tp}[chapter]
\renewcommand{\thefp}{\thechapter.\arabic{fp}}
\renewcommand{\thetp}{\thechapter.\arabic{tp}}

\newtcolorbox{finepoint}[1][]{
  enhanced, breakable, colback=fpbg, colframe=fpbg,
  colbacktitle=fpbg, coltitle=black,
  boxrule=0pt, arc=0pt, left=9pt, right=9pt, top=5pt, bottom=8pt,
  toptitle=3pt, bottomtitle=3pt,
  before skip=10pt, after skip=10pt,
  fontupper=\small, fonttitle=\heiti\small,
  title={\refstepcounter{fp}\FPname~\thefp\ifx\relax#1\relax\else\hspace{0.8em}#1\fi},
}

\newtcolorbox{techpoint}[1][]{
  enhanced, breakable, colback=tpbg, colframe=rule!70,
  colbacktitle=tpbg, coltitle=black,
  boxrule=0.5pt, arc=0pt, left=9pt, right=9pt, top=5pt, bottom=7pt,
  toptitle=3pt, bottomtitle=3pt,
  before skip=10pt, after skip=10pt,
  fontupper=\small, fonttitle=\heiti\small,
  title={\refstepcounter{tp}\TPname~\thetp\ifx\relax#1\relax\else\hspace{0.8em}#1\fi},
}

% 定义小方框（原书置于边栏的定义式）
\newtcolorbox{defbox}[1][]{
  enhanced, colback=white, colframe=black!55,
  boxrule=0.5pt, arc=0pt, left=6pt, right=6pt, top=4pt, bottom=4pt,
  halign=center, fontupper=\small,
  before skip=8pt, after skip=8pt,
  title={#1}, coltitle=black, fonttitle=\small\heiti,
  attach title to upper={\par\vspace{2pt}}
}

% ---------- 数学算子 ----------
\newcommand{\dox}[1]{\mathrm{do}(#1)}
\newcommand{\pr}{\Pr}
\newcommand{\indep}{\perp\!\!\!\perp}
\newcommand{\nindep}{\not\perp\!\!\!\perp}
\newcommand{\E}{\mathrm{E}}
\newcommand{\Var}{\mathrm{Var}}
\newcommand{\given}{\,|\,}

% ---------- TikZ 统一风格 ----------
\tikzset{
  causal/.style={>={Stealth[length=2.1mm,width=1.5mm]},semithick},
  vtx/.style={circle,draw,inner sep=1pt,minimum size=6.2mm,font=\small},
  vtxs/.style={circle,draw,inner sep=1pt,minimum size=5.2mm,font=\footnotesize},
  lat/.style={circle,draw,dashed,inner sep=1pt,minimum size=6.2mm,font=\small},
  ed/.style={->,semithick},
  edb/.style={->,very thick},
  edd/.style={->,semithick,dashed},
  bi/.style={<->,semithick},
  bid/.style={<->,semithick,dashed},
}

% ---------- 杂项 ----------
\renewcommand{\arraystretch}{1.12}
\setlist[itemize]{leftmargin=1.6em,itemsep=0.15em,topsep=0.3em}
\setlist[enumerate]{leftmargin=1.9em,itemsep=0.15em,topsep=0.3em}

% ---------- 超链接（最后）----------
\usepackage{url}
\usepackage[bookmarks=true,bookmarksnumbered=true,colorlinks=true,
            linkcolor=inkblue,citecolor=inkblue,urlcolor=inkblue]{hyperref}

% 定义框内禁止数学间距被拉伸
\AtBeginEnvironment{defbox}{\medmuskip=4mu \thickmuskip=5mu}
